% Supporting checking notes; read alongside paper/paper.tex.
\section{Computation and formalization}\label{sec:formal}
This section specifies which parts are machine-checked and what the
checks assume. The mathematical upper bound is formalized from the
adjacency matrix and spectral definition of strong cospectrality;
the attaining example is a separate exact computation.
The complete accompanying bundle contains the paper, the Lean project,
the certificate data and the exact checkers. All paths in this paper
are relative to that bundle: the formal definitions and theorem are
in \texttt{lean/CubelikeSC/Statement.lean} and
\texttt{lean/CubelikeSC/Proof.lean}, the certificates are in
\texttt{data/d7k3.cert}, and the arithmetic check is
\texttt{checks/check.py}. The bundle's \texttt{README.md} gives the
commands for checking it as a standalone directory.

\subsection{The formal upper-bound statement}
The formal definitions identify the machine-checked statement with
Theorem~\ref{thm:main}'s upper bound. The project uses Lean~4 and
mathlib, with versions fixed by its Lake configuration. Its definitions
are reproduced below, with line breaks for readability:
\begin{lstlisting}
abbrev V := Fin 7 → ZMod 2

def adj (S : Finset V) : Matrix V V ℝ :=
  fun x y => if x - y ∈ S then 1 else 0

def StronglyCospectral (S : Finset V) (u v : V) : Prop :=
  ∀ θ : ℝ,
    Module.End.HasEigenvalue (Matrix.toEuclideanLin (adj S)) θ →
    let E := Module.End.eigenspace (Matrix.toEuclideanLin (adj S)) θ
    E.starProjection (EuclideanSpace.single u 1) =
      E.starProjection (EuclideanSpace.single v 1) ∨
    E.starProjection (EuclideanSpace.single u 1) =
      -E.starProjection (EuclideanSpace.single v 1)

def EndpointSC : Prop :=
  ∀ S : Finset V, (0 : V) ∉ S → ∀ P : Finset V,
    (∀ u ∈ P, ∀ v ∈ P, StronglyCospectral S u v) → P.card ≤ 4
\end{lstlisting}
The adjacency convention $x-y\in S$ agrees with $x+y\in S$ over
$\F_2$. The expression \texttt{E.starProjection} is the orthogonal
projection onto the real eigenspace \texttt{E}; the disjunction permits
either sign separately for each eigenvalue.
The name \texttt{EndpointSC} denotes the final universal upper-bound
assertion. Its proof is declared as
\begin{verbatim}
theorem endpoint_sc : EndpointSC
\end{verbatim}
The definition of \texttt{EndpointSC} is the universal statement
\[
 \forall S\subseteq V,\quad 0\notin S\ \Longrightarrow\
 \forall P\subseteq V,\quad
 (\forall u,v\in P,\ u,v\text{ are strongly cospectral})
 \ \Longrightarrow\ |P|\leq4.
\]
It quantifies over every connection set excluding zero and every
pairwise strongly cospectral set, including disconnected graphs.

The formalization proves the direction of
Lemma~\ref{lem:criterion} needed for the upper bound:
strong cospectrality of $u,v$ implies $v-u\in K(S)$.
It proves the subspace and parity facts, the moment identities, bracket
extraction, the penalty bound, the finite certificate checks, and their
case cover. Instead of changing basis, it selects three independent
vectors of $K(S)$ whose span contains $\sigma$, and indexes character
cosets by their three dot products. This is the same reduction in a
coordinate-independent form.

All certificate calculations are evaluated by Lean's kernel. No axiom asserts the finite inequality, and
\texttt{native\_decide} is not used. The final theorem depends only on the usual
logical axioms \texttt{propext}, \texttt{Classical.choice}, and
\texttt{Quot.sound}. The build fails if the theorem's axiom dependencies
differ from this list.
From \texttt{lean/}, the standard check is
\begin{verbatim}
lake exe cache get
lake build
\end{verbatim}
The trusted foundation for that check is Lean's kernel and the
definitions and library results used by the proof. The formal theorem's
identification with the graph-theoretic statement is given by the
definitions just described, rather than by a hypothesis about character
eigenvalue fibers.

\subsection{Independent arithmetic checks}
The separate integer checker verifies the certificates directly for the
original cases, and also checks the rejection of deliberately damaged data.
The resulting certificate counts are summarized in Table~\ref{tab:certificates}.
\begin{table}[htbp]
\centering
\begin{tabular}{lr}
\hline
Certificate disposition & Number of case keys\\
\hline
A coset has no allowed bracket & $58$\\
Nonnegative weights satisfy~\eqref{eq:scaleddual} & $3{,}926$\\
\hline
Total & $3{,}984$\\
\hline
\end{tabular}
\caption{Exact certificate outcomes. Each key can represent several
cases; these are counts of keys, not counts of connection sets or cases.}
\label{tab:certificates}
\end{table}
The $58$ empty-list entries have $n=0,1,2,3,4,5,6$, with respectively
$2,5,9,15,16,9,2$ keys. Together they represent $1{,}628$ original cases.
For each entry, forming the relevant $J_u$ verifies that no integer pair
satisfies conditions 1--4. For example, at $n=0$ the only possible endpoint
is zero, which condition~4 forbids on a nonzero coset.
A separate exact Python checker forms each case's own eight bracket
lists and checks the matching certificate on those lists, without
assuming the key-sharing argument.

\subsection{Certificate generation and the lower example}
We distinguish the untrusted generation of certificates from their
exact verification, and record the separate check of attainment.
The certificate weights were obtained from a floating-point linear
programming dual and converted to rational numbers.
This generation step is not trusted: only the inequalities on the
supplied integer weights are used. The separate Python certificate
checker uses arbitrary-precision integers. Its role is an additional
implementation of the certificate calculation, not an additional
mathematical assumption in the Lean theorem.
To test rejection, the supplied checks deliberately remove an entry,
falsely assert that a bracket list is empty, weaken or zero weights,
make a weight negative, and duplicate an entry. Every altered file is
rejected. Lean also proves that selected entries fail their check when
weights are removed or halved, and that deleting a required entry
leaves a case without a certificate. The command
\texttt{python3 checks/check.py} performs the exact Python checks
and these rejection tests; \texttt{checks/README.md} explains them.

Earlier computational work is recorded in \'Arnad\'ottir and Godsil's
Appendix~A, which gives six connection sets on $32$ vertices; their
complements give the other six graphs. Appendix~B provides Sage code.
Their acknowledgments credit Gordon Royle's cubelike graph lists.

The converse of Lemma~\ref{lem:criterion} and attainment are not part
of the Lean development. For the displayed construction, exact
character sums check~\eqref{eq:attain}; a separate computation checks
strong cospectrality using spectral projection columns with rational
entries. The construction and the written character argument above
give the lower bound four. Neither the upper theorem nor this lower
example relies on a numerical eigensolver.

Figure~\ref{fig:example} and its SVG version in \texttt{data/figure.svg}
come from the same graph data. The optional illustration script
\texttt{checks/figure.py} checks the graph's vertex count, edge count,
degree, connectivity and highlighted class, and reproduces the figure
and the spectral-sign table. It is not needed for the upper-bound proof.


\section{Further presentation notes}
Condition~2 is
redundant given condition~1; it is retained to match the finite assertion
checked by the accompanying certificates.

Cases with no admissible assignment satisfy it vacuously.
The parameter range includes $n=0$, $B=\varnothing$, $t=0$, and all
other boundary values; all of them are included in the finite check.



The data are supplied in
\texttt{data/d7k3.cert} in the accompanying bundle.


\section{Verification and reproducibility}\label{sec:formal}
Lean and mathlib prove the universal upper bound, with formal theorem
\texttt{CubelikeSC.endpoint\_sc}:
\begin{verbatim}
theorem endpoint_sc : EndpointSC
\end{verbatim}
The definition of \texttt{EndpointSC} is the universal statement
\[
 \forall S\subseteq V,\quad 0\notin S\ \Longrightarrow\
 \forall P\subseteq V,\quad
 (\forall u,v\in P,\ u,v\text{ are strongly cospectral})
 \ \Longrightarrow\ |P|\leq4.
\]
Here the formal vertex type is $(\mathbb Z/2\mathbb Z)^7$.
The adjacency convention $x-y\in S$ agrees with $x+y\in S$ over
$\F_2$. Strong cospectrality is defined using orthogonal projections
onto the real adjacency eigenspaces, with either sign permitted for
each eigenvalue. Thus the formal assertion is precisely the upper bound
of Theorem~\ref{thm:main}, with no connectivity hypothesis.

The formalization proves the direction of
Lemma~\ref{lem:criterion} needed for the upper bound:
strong cospectrality of $u,v$ implies $v-u\in K(S)$.
It proves the subspace and parity facts, the moment identities, bracket
extraction, the penalty bound, the finite certificate checks, and their
case cover. Instead of changing basis, it selects three independent
vectors of $K(S)$ whose span contains $\sigma$, and indexes character
cosets by their three dot products. This is the same reduction in a
coordinate-independent form.

The converse character criterion and attainment remain outside Lean.
The written character argument and \'Arnad\'ottir and Godsil's construction
give the lower bound; independent exact computations check the character
sums and the adjacency spectral projections. A separate integer checker
verifies every original case's certificate and rejects damaged data.
The checking guide gives the reproduction commands, formal definitions,
logical dependencies and certificate details. The proof uses Lean's
kernel and its library results; no solver output is assumed as an axiom.



\section{Historical and dimensional context}
The question appeared in their first preprint, posted September 15, 2021;
the journal article appeared in 2023.
This result does not bound the size of strongly cospectral sets uniformly
over all dimensions. The maximum in dimensions eight and higher is not determined
here.
